%%%%%%%%%%%%%%%%%%%%%%%%%%%%%%%%%%%%%%%%%%%%%%%%%%%%%%%%%%%%%%%%%%%%%%%%%%%%%%
%
% 02_introduction.tex
%
%%%%%%%%%%%%%%%%%%%%%%%%%%%%%%%%%%%%%%%%%%%%%%%%%%%%%%%%%%%%%%%%%%%%%%%%%%%%%%

\begin{mosbox}{1. Motivation: Why Represent Memory Geometrically?}

\BodyFont

Modern learning systems typically represent memory as collections of
vectors or key--value databases. Such representations are highly effective
for retrieval but contain little explicit mathematical structure governing
how individual concepts relate, how local consistency propagates, or how
global coherence can be verified.

MOS investigates a different mathematical viewpoint.

Rather than treating memory as an unordered collection of embeddings,
memory is modeled as a structured geometric object

\[
M=(C,\mathcal F,s),
\]

where relationships between concepts become intrinsic features of the
space itself.

\vspace{4mm}

The objective is not to replace existing reasoning systems.
Instead, MOS proposes an external mathematical organization layer capable
of

\begin{itemize}

\item organizing knowledge hierarchically,

\item measuring consistency,

\item identifying contradictions,

\item supporting principled memory growth,

\item providing mathematically interpretable diagnostics.

\end{itemize}

\vspace{4mm}

This perspective shifts the emphasis

\[
\boxed{
\text{Memory}
\longrightarrow
\text{Structured Mathematical Object}
}
\]

rather than

\[
\boxed{
\text{Memory}
\longrightarrow
\text{Collection of Independent Vectors}.
}
\]

\end{mosbox}

%%%%%%%%%%%%%%%%%%%%%%%%%%%%%%%%%%%%%%%%%%%%%%%%%%%%%%%%%%%%%%%%%%%%%%%%%%%%%%
%% GEOMETRIC VIEW
%%%%%%%%%%%%%%%%%%%%%%%%%%%%%%%%%%%%%%%%%%%%%%%%%%%%%%%%%%%%%%%%%%%%%%%%%%%%%%

\vspace{5mm}

\begin{center}

\begin{tikzpicture}[scale=1.15]

%%%%%%%%%%%%%%%%%%%%%%%%%%%%%%%%%%%%%%%%%%%%%%%%%%%%%%%%%%%%
%% VECTOR DATABASE
%%%%%%%%%%%%%%%%%%%%%%%%%%%%%%%%%%%%%%%%%%%%%%%%%%%%%%%%%%%%

\node at (-7,5)
{\Large\bfseries Conventional Memory};

\foreach \i in {0,...,13}
{
\fill[MOSOrange]
(-8+rand*2,3+rand*2)
circle(2.5pt);
}

\draw[very thick,MOSRed]
(-9,2)
rectangle
(-5,6);

\node at (-7,1.2)
{\Large Independent Embeddings};

%%%%%%%%%%%%%%%%%%%%%%%%%%%%%%%%%%%%%%%%%%%%%%%%%%%%%%%%%%%%
%% ARROW
%%%%%%%%%%%%%%%%%%%%%%%%%%%%%%%%%%%%%%%%%%%%%%%%%%%%%%%%%%%%

\draw[arrow]
(-4,4)
--
(0,4);

\node at (-2,4.6)
{\Large Structure};

%%%%%%%%%%%%%%%%%%%%%%%%%%%%%%%%%%%%%%%%%%%%%%%%%%%%%%%%%%%%
%% MOS
%%%%%%%%%%%%%%%%%%%%%%%%%%%%%%%%%%%%%%%%%%%%%%%%%%%%%%%%%%%%

\node at (6,5)
{\Large\bfseries MOS Representation};

\coordinate (A) at (4,2.5);

\coordinate (B) at (8,2.5);

\coordinate (C) at (6,5.5);

\coordinate (D) at (6,3.8);

\draw[simplex]
(A)--(B)--(C)--cycle;

\draw[simplex]
(A)--(D);

\draw[simplex]
(B)--(D);

\draw[simplex]
(C)--(D);

\foreach \p in {A,B,C,D}
{

\fill[MOSOrange]
(\p)
circle(3.5pt);

}

\node at (6,1.2)
{\Large Higher-Order Relations};

\end{tikzpicture}

\end{center}

%%%%%%%%%%%%%%%%%%%%%%%%%%%%%%%%%%%%%%%%%%%%%%%%%%%%%%%%%%%%%%%%%%%%%%%%%%%%%%
%% REMARK
%%%%%%%%%%%%%%%%%%%%%%%%%%%%%%%%%%%%%%%%%%%%%%%%%%%%%%%%%%%%%%%%%%%%%%%%%%%%%%

\begin{remarkbox}

\BodyFont

\textbf{Interpretation}

The role of the simplicial complex is not merely to connect concepts by
edges, but to encode higher-order interactions involving multiple concepts
simultaneously. The cellular sheaf then equips each cell with local data,
while global sections describe mutually compatible knowledge across the
entire structure.

\end{remarkbox}